\documentclass[11pt,a4paper]{article}
\usepackage[utf8]{inputenc}
\usepackage[T1]{fontenc}
\usepackage[margin=22mm]{geometry}
\usepackage{lmodern,xcolor,enumitem}
\usepackage[colorlinks=true,urlcolor=teal,linkcolor=teal,pdfborder={0 0 0}]{hyperref}
\definecolor{ink}{HTML}{173A3A}
\definecolor{teal}{HTML}{006A70}
\setlength{\parindent}{0pt}
\setlength{\parskip}{7pt}
\setlist[enumerate]{leftmargin=*,itemsep=4pt,topsep=3pt}
\setlist[itemize]{leftmargin=*,itemsep=3pt,topsep=3pt}
\hypersetup{pdftitle={Program Matematika Indonesia: Panduan belajar / Learning guide},
 pdfauthor={Program Matematika Indonesia},pdfsubject={Panduan peserta belajar dan pengajar; learner and teacher guide}}
\pdfinfoomitdate=1
\pdftrailerid{}
\pdfsuppressptexinfo=15
\newcommand{\heading}[1]{\par\medskip{\large\bfseries\color{ink}#1}\par\smallskip}
\newcommand{\idsite}{https://kokunoyumeto.github.io/program-matematika-indonesia/id/}
\newcommand{\ensite}{https://kokunoyumeto.github.io/program-matematika-indonesia/en/}
\newcommand{\backsite}{https://kokunoyumeto.github.io/program-matematika-indonesia/backend/}
\newcommand{\doiurl}{https://doi.org/10.5281/zenodo.22059707}
\begin{document}
\raggedright
{\small\color{teal}PANDUAN BELAJAR \hfill VERSI 0.63.32 / 30 SEPTEMBER 2026}

{\LARGE\bfseries\color{ink}Program Matematika Indonesia}

Jalur belajar terbuka, dari dasar hingga matematika lanjut.
Pilih bahan sesuai pengetahuan awal dan tujuan Anda; tidak perlu mengikuti
semua mata kuliah dalam satu urutan.

\href{\idsite}{\textbf{Mulai dalam Bahasa Indonesia}}
\quad | \quad \href{\ensite}{English}
\quad | \quad \href{\doiurl}{Arsip unduhan}

\heading{Mulai belajar}
\begin{enumerate}
\item \textbf{Pilih mata kuliah.} Buka kartu mata kuliah atau peta belajar.
Periksa bagian prasyarat untuk menemukan pengetahuan yang dipakai.
\item \textbf{Buka bahan bacaan.} Gunakan tautan pembaca daring atau unduhan
yang tercantum. PDF, HTML, EPUB dan sumber yang dapat disunting tidak tersedia
dalam bentuk yang sama untuk setiap buku. Bahasa bahan bacaan ditandai terpisah
dari bahasa antarmuka.
\item \textbf{Kerjakan latihan.} Gunakan latihan dan alat belajar yang tersedia
untuk mata kuliah tersebut. Petunjuk, jawaban dan solusi bukan hal yang sama.
Ruang jawaban kosong tidak berarti solusi telah disediakan.
\item \textbf{Tentukan langkah berikutnya.} Kembali ke kartu mata kuliah untuk
memilih topik lanjutan. Catatan kemajuan pribadi tersimpan di peramban, bukan
di server; catatan itu bukan penilaian atau sertifikat.
\end{enumerate}

\heading{Untuk pengajar dan pengguna ulang}
\href{\backsite}{Indeks backend bersama} memuat alat peserta belajar dan bahan
pengajar yang telah dihubungkan. Pada mata kuliah yang menyediakannya,
perencana latihan mempertahankan identitas soal serta status petunjuk dan
solusi dari sumber. Periksa batas cakupan pada setiap alat sebelum memakainya.
Sumber asli, versi, lisensi dan catatan koreksi tetap melekat pada tiap bahan;
program ini tidak menggantinya dengan satu lisensi untuk semua buku.

\heading{Luring, sumber, dan batas yang masih ada}
Peta belajar satu berkas dapat diunduh dari halaman program. Paket navigator
portabel memiliki \texttt{START-HERE.html}. Paket backend bersama berbeda:
ekstrak ZIP, jalankan \texttt{python -m http.server 8000 -{}-directory docs},
lalu buka \texttt{http://localhost:8000/backend/}. Cara ini memerlukan Python
yang telah terpasang dan menyediakan data JSON bagi alat interaktif.

Buku yang tidak ada di dalam paket harus diunduh terpisah. Tautan luar dan
penampil rumus D50 yang masih memakai MathJax dari CDN memerlukan jaringan.
ZIP backend menyertakan panduan pembangunan ulang dan identitas berkas;
bukan salinan lengkap seluruh korpus buku.

Saat rilis ini dibuat, \textbf{40 peran mata kuliah} memiliki adaptor bersama
yang tervalidasi. Kesetaraan kemampuan backend asli baru tercatat untuk
\textbf{15 dari 40 peran}. Ini bukan persentase penerjemahan, bukan pernyataan
bahwa semua backend selesai, dan bukan bukti peninjauan kebahasaan oleh ahli.

\vfill
{\footnotesize Panduan ini disusun oleh OpenAI Codex --- gpt-6-astra, Ultra effort.
Tidak diklaim ada penyuntingan manusia. Kredit dan riwayat penerjemahan setiap
buku tetap mengikuti catatan edisinya. Sumber LaTeX panduan tersedia langsung
di samping PDF; ZIP sumber lengkap memuat berkas dan petunjuk pembangunannya.}

\newpage
{\small\color{teal}LEARNING GUIDE \hfill VERSION 0.63.32 / 30 SEPTEMBER 2026}

{\LARGE\bfseries\color{ink}An open mathematics program}

A route from foundations to advanced mathematics. Choose materials for your
background and aims; you do not need to take every course in a single sequence.

\href{\ensite}{\textbf{Start in English}}
\quad | \quad \href{\idsite}{Bahasa Indonesia}
\quad | \quad \href{\doiurl}{Download archive}

\heading{Start learning}
\begin{enumerate}
\item \textbf{Choose a course.} Open a course card or the learning map.
Check the prerequisites to see which earlier ideas the course uses.
\item \textbf{Open the reading.} Use the listed online reader or download.
PDF, HTML, EPUB and editable source are not available in the same combination
for every book. A reading's language is identified separately from the
interface language.
\item \textbf{Try the exercises.} Use the exercises and learning tools provided
for that course. Hints, answers and worked solutions are different things.
An empty response box does not mean a solution has been supplied.
\item \textbf{Choose your next step.} Return to the course card for onward
topics. Personal progress stays in your browser, not on a server; it is not
an assessment or a certificate.
\end{enumerate}

\heading{For teachers and people reusing the materials}
The \href{\backsite}{shared backend index} links the integrated learner tools
and teaching materials. Where a course provides them, exercise planners retain
source question IDs and distinguish supplied hints and solutions from missing
ones. Check each tool's scope notes before using it.

Original sources, editions, licences and correction records stay attached to
each resource. The program does not replace them with a blanket licence for
all books. An English interface does not imply that every linked resource has
an English edition.

\heading{Offline use, source files, and remaining limits}
Download the single-file learning map from the program page. The portable
navigator package has a \texttt{START-HERE.html}. The shared backend package
is different: extract its ZIP, run
\texttt{python -m http.server 8000 -{}-directory docs}, then open
\texttt{http://localhost:8000/backend/}. This requires Python to be installed
and makes JSON data available to the interactive tools.

Books outside the package need separate downloads. External links and D50's
MathJax CDN still need an internet connection. The backend ZIP includes replay
instructions and file identities, not a complete copy of all the books.

At this release, \textbf{all 40 course roles} have validated common adapters.
Parity with the original backend capabilities is recorded for
\textbf{15 of 40 roles}. This is not a translation percentage, a claim that
every backend is finished, or evidence of expert language review.

\vfill
{\footnotesize This guide was prepared by OpenAI Codex --- gpt-6-astra, Ultra effort.
No human editing is claimed. Each book retains its own translation and contributor
history. The guide's complete LaTeX is directly available beside the PDF; its
source ZIP supplies the source and build instructions.}
\end{document}
